\section{Abstract}

Movie recommendation systems must predict user preferences from sparse, long-tailed interaction data.
This project presents \textbf{MovieMind}, an end-to-end pipeline on the MovieLens 1M dataset that combines
exploratory analysis, leakage-safe feature engineering, classical machine learning, and neural collaborative
filtering (NCF).
We evaluate models with root mean squared error (RMSE) and mean absolute error (MAE) on a chronological
70/10/20 train/validation/test split.
Among all families tested, \textbf{raw Gradient Boosting} on engineered tabular features achieves the best
test performance (RMSE \textbf{0.8981}, MAE \textbf{0.7062}).
Hyperparameter tuning improves the deep learning path but does not surpass the tree-based champion.
The full workflow is reproduced in \texttt{notebooks/MovieMind\_capstone.ipynb}; an optional FastAPI and
Streamlit deployment extends the work beyond offline modeling.

\textbf{Keywords:} recommender systems, MovieLens, gradient boosting, neural collaborative filtering, cold start, long tail.
